%-------------------------------------------------------------------------------
\cvsection{工作经历}

%---------------------------------------------------------
\begin{cventries}

%---------------------------------------------------------
  \cventry
    {AI 应用 \& 平台架构师（技术负责人）} % Job title
    {垚阳培极科技（深圳）} % Organization
    {深圳 · 福田} % Location
    {03. 2021 - 至今} % Date(s)
    {
      \begin{cvitems}
        \item {负责公司 AI 应用与平台技术架构设计，主导「端侧智能体平台 + 业务中台 + DevOps 体系」三条主线，统筹需求评审、架构设计、核心编码与技术文档全流程。}
        \item {AI 方向：基于 Spring Boot / Python 构建大模型应用服务与 Agent 编排能力，落地意图识别、对话记忆（RAG + 向量检索）、音色/角色切换与多轮对话管理；主导端云协同方案设计，对接 ESP32 端侧智能硬件。}
        \item {平台方向：设计并二开企业级商城 / OA 微服务中台（产品、订单、支付、权限、营销等 20+ 模块），采用前后端分离 + 微服务架构 + Seata 分布式事务 + Redis 分布式缓存，支撑 250W 级行业合同交付。}
        \item {工程效能：建设自动化 CI/CD 体系（代码托管、静态扫描、自动化测试、容器化部署），Docker + Nginx 实现标准化交付，显著降低发布耗时与线上故障率。}
      \end{cvitems}
    }

%---------------------------------------------------------
  \vspace{3mm}
  \cventry
    {高级软件工程师（算法与平台组）} % Job title
    {雅乐美森（服务方：天津大学 / 人民银行数据中心 / 金融科技机构）} % Organization
    {北京 · 中关村} % Location
    {09. 2017 - 12. 2020} % Date(s)
    {
      \begin{cvitems}
        \item {金融科技：主导高并发交易转账模块研发，基于 Redis 分布式锁 + Seata 分布式事务 + Crontab 对账，达成 \textbf{QPS 10W+、TPS 1W+}，数据误差容忍度 0.01；Jmeter 压测与全链路性能优化。}
        \item {算法与量化：基于 Python / Pandas / NumPy / Scikit-learn / TensorFlow 完成数据处理、特征工程、机器学习模型训练评估及期货量化策略回归测试。}
        \item {舆情与风控 NLP：负责数据采集（Scrapy）、文本处理（Goose3 / NLTK）、情感分析（VADER / SnowNLP）、主题建模（BERTopic）与知识图谱（Neo4j）分析，Kibana / Grafana 可视化落地。}
      \end{cvitems}
    }

%---------------------------------------------------------
\end{cventries}
\vspace{3mm}
